\section{Implementation}
\label{sec:implementation}

The implementation has six load-bearing components: a runtime kernel that exposes a pre-dispatch admission hook and a treaty primitive; a federation layer that performs strict bilateral DSSE verification; a selective-disclosure module over a BBS message vector; a multi-lane anchor for public-witness evidence; a fixture-backed three-vendor case-study generator; and the Lean proof root. Each implementation citation below resolves to a \codepath{file:line}, and each citation falls into one of three categories: a runtime obligation that discharges a Lean theorem, a tested executable path, or a fixture. Table~\ref{tab:artifacts} indexes the artifacts; Table~\ref{tab:theorems} maps them to Lean theorems.

\begin{table*}[t]
\caption{Load-bearing implementation artifacts used by the paper.}
\label{tab:artifacts}
\scriptsize
\begin{tabularx}{\textwidth}{@{}>{\raggedright\arraybackslash}p{.24\textwidth}>{\raggedright\arraybackslash}p{.38\textwidth}>{\raggedright\arraybackslash}X@{}}
\toprule
Artifact & Enforced statement & Evidence boundary \\
\midrule
\codepath{TreatyScope} & Participants, public keys, ladder-manifest hashes, allowed action classes, validity window, revocation epoch, and trust bundle are bound before intersection. & Schema and semantic hash checks are in the runtime; treaty law outside these fields is not represented. \\
\codepath{LadderIntersection} & For each allowed action class, the joint mode, consistency model, co-sign rule, evidence set, and per-participant mode are materialized. & Correct only for the finite ladder schema accepted by the current validators. \\
\codepath{CrossKernelContinuation} & Cross-kernel child work is linked to parent receipt, session anchor, capability id, action class, target tool, nonce, and validity window. & Prevents replay only when the receiving runtime consumes or releases continuation state correctly. \\
\codepath{BilateralInvocation} & Local and remote receipts, request and outcome hashes, lineage statement, continuation, action class, consistency model, and signer ids are jointly bound. & Verifies joint intent; it does not prove the underlying human commercial contract. \\
\codepath{SelectiveDisclosureProof} & A verifier can check disclosed projection fields against an issuer key without receiving every receipt field. & Privacy proof is secondary to the authoritative Ed25519 receipt. \\
\codepath{AnchorWitnessPolicy} & Public witness lanes must verify active inclusion or use explicitly cached verifier-owned evidence. & Rekor SET verification is stronger than OTS advisory mode; lane-specific limitations remain. \\
\bottomrule
\end{tabularx}
\end{table*}

\paragraph{Treaty primitive.}
The treaty module \codepath{crates/kernel/chio-runtime-core/src/treaty.rs} discharges the runtime side of \thm{treaty_admission_iff_predicate_intersection}. It hashes treaty scope with canonical JSON, computes semantic hashes that exclude validity windows for stable intersection identity, and binds bilateral invocations over schema, invocation id, treaty id, ladder-intersection hash, continuation hash, action class, consistency model, capability id, request hash, outcome hash, local and remote receipt hashes, and signer ids (\codepath{crates/kernel/chio-runtime-core/src/treaty.rs:81}). \codepath{validate_treaty_scope} rejects unsupported schemas, invalid windows, missing participants, duplicate participants, hash mismatches, and malformed trust bundle references (\codepath{crates/kernel/chio-runtime-core/src/treaty.rs:188}). \codepath{compute_ladder_intersection} then checks participant manifests and constructs the joint action-class table.

The treaty module turns political metaphors into ordinary data validation. A treaty with one participant rejects; with duplicate participant keys, rejects; with a governance ladder hash unequal to the canonical hash of the manifest, rejects; with an action class absent from one participant's manifest, rejects; with a destructive action whose joined mode falls below receipt-backed, rejects. A programmable constitution fails as a schema-and-predicate matter before it becomes a governance debate.

The semantic-hash split separates stable identity from liveness. Treaty-scope semantic identity excludes the validity window so that short-lived renewals carry the same policy identity while admission still enforces freshness. Stable identity belongs to the predicate set; liveness belongs to the admission context.

\paragraph{Admission hook.}
\codepath{crates/kernel/chio-runtime-core/src/admission_hook.rs} implements \codepath{RuntimeAdmissionHook}. It permits non-Chio requests only when no Chio runtime context is present; otherwise malformed context denies. Malformed treaty context denies with \codepath{invalid_chio_treaty_context} or a missing-field code such as \codepath{missing_treaty_scope_id}. Request-level smuggling of trust roots, peer directories, signing keys, or dynamic trust bundles is rejected by \codepath{treaty_ref_from_request} (\codepath{crates/kernel/chio-runtime-core/src/admission_hook/treaty_ref.rs:24}). Treaty DSSE evidence is verified against verifier-owned store artifacts; signer sets must equal treaty participants and public keys must be independent (\codepath{crates/kernel/chio-runtime-core/src/admission_hook/dsse.rs:87}). The hook releases reserved continuation state after runtime denial or abort, closing a replay edge tested in \codepath{crates/kernel/chio-runtime-core/tests/runtime_admission.rs}.

The hook is the point at which the model is falsifiable against real code. Allowing federated origin without treaty context would make the treaty-intersection theorem irrelevant to cross-kernel execution. Allowing request metadata to smuggle a trust root would collapse the verifier-owned-store assumption. Failing to release reserved continuation state on denial would create a state leak and ambiguous replay behavior. Each of these is a concrete place where agent middleware turns formal policy into advisory logging.

The hook is asymmetric: non-Chio local requests use the old path when no Chio context is present, and once Chio context appears, malformed or incomplete context denies. Backward compatibility is preserved; partial federation is not a bypass. The receiver's kernel decides when a request has crossed into the constitutional path.

\paragraph{Strict DSSE verifier.}
\codepath{crates/trust/chio-federation/src/bilateral_dsse.rs} defines the strict predicate and envelope used by every cross-organization invocation. \codepath{verify_chio_bilateral_dsse_envelope} rejects wrong payload type, wrong signature count, noncanonical JSON, wrong statement type, wrong predicate type, missing subject, and reused signer keys (\codepath{crates/trust/chio-federation/src/bilateral_dsse.rs:1167}). \codepath{crates/trust/chio-federation/src/bilateral_verifier.rs} layers operational checks on top: peer pins, fresh ladder manifests, revocation epoch, receipt resolution and signature, tool argument hash, policy verdict agreement, capability lease, governance receipt for receipt-backed classes, and treaty binding refs (\codepath{crates/trust/chio-federation/src/bilateral_verifier.rs:872}). Together, the envelope check and the operational check ensure that admission depends on the predicate the protocol names, not on a strict subset that a sender could satisfy with degenerate data.

The strict verifier refuses ``signed JSON'' as an adequate security boundary. A DSSE envelope with the wrong predicate is not a near miss; a payload signed by the same key twice is not two-party consent; a receipt with a valid signature but the wrong request hash is not admissible under the treaty; receipt-backed classes require a governance receipt even when the ordinary receipt is otherwise valid. The error taxonomy is part of the protocol because downstream audit must know which constitutional precondition failed.

\paragraph{Buyer audit and three-vendor loop.}
The three-vendor example is a generator in \codepath{examples/chio-3vendor/src/main.rs} that emits buyer-auditor proof packages, workflow-intersection JSON, signed negative packages, pheromone artifacts, and runtime-spine fixtures. The example is gated by integration tests that re-run the closure on every change to any of its bound components.

Buyer closure is the concrete example of scoped sovereignty. The buyer demands the specific evidence package required for admission, not every internal receipt; the package binds local and remote receipts to treaty scope, ladder intersection, continuation, and lineage. A system that cannot generate signed denials cannot prove it rejects stale, mismatched, or replayed evidence; the generator produces both positive and negative packages in the same canonical schema.

The example demonstrates that the runtime and federation crates agree on artifact shape, emit repeatable fixtures, and exercise negative cases. Because every bound component (treaty scope, ladder intersection, bilateral DSSE, admission hook, buyer audit) is exercised, a regression in any one breaks the integration gate.

\paragraph{Proof inventory.}
The Lean proof root contains the theorem and lemma declarations the runtime depends on, including the ones introduced for this work, with no admitted goals (\codepath{sorry}) and no extra axioms (\codepath{axiom}). A curated inventory bridges formal names to protocol properties and Rust symbols; the curated subset is the audited coverage the paper defends.

Cryptographic primitives, canonical-JSON implementation correctness, TLS, SQLite, public-chain finality, and several other lower layers are outside the Lean proof root (Table~\ref{tab:assumptions}). Selected constitutional predicates and protocol invariants have named Lean theorems, no unchecked placeholders, and inventory entries linking those theorems to protocol properties.

\begin{table*}[t]
\caption{Selected proof-to-code map. The first four rows are theorems introduced in this work.}
\label{tab:theorems}
\scriptsize
\begin{tabularx}{\textwidth}{@{}>{\raggedright\arraybackslash}p{.28\textwidth}>{\raggedright\arraybackslash}p{.35\textwidth}>{\raggedright\arraybackslash}X@{}}
\toprule
Theorem & Informal statement & Code or property \\
\midrule
\thm{treaty_admission_iff_predicate_intersection} & Treaty admission equals treaty, left-polity, and right-polity predicate intersection. & \codepath{crates/kernel/chio-runtime-core/src/treaty.rs:261} \\
\thm{treaty_admission_stable_under_ladder_floor} & Satisfied ladder floor preserves treaty admission. & \codepath{crates/kernel/chio-runtime-core/src/treaty.rs:410} (\codepath{validate_ladder_intersection}) \\
\thm{amendment_admissible_iff_backward_refinement} & Amendment admission iff new predicates imply old predicates. & Polity amendment model \\
\thm{amendment_without_refinement_rejected} & Missing proof rejects enactment. & Crisis admission invariant \\
\thm{proof.delegation_step_allow_requires_attenuation} & Delegation allow implies attenuation. & Capability delegation \\
\thm{theorem.budget.sibling_sum_soundness} & Child budget sum stays within parent share. & \codepath{BudgetSplit::try_admit_child} \\
\thm{proof.receipt_sign_then_verify} & Signed symbolic receipt verifies. & Receipt signing \\
\thm{proof.verified_receipt_lineage_sound} & Verified lineage requires both receipts and linkage. & Buyer lineage audit \\
\bottomrule
\end{tabularx}
\end{table*}
